\documentclass[a4paper,12pt]{article}
\usepackage[utf8]{inputenc}
\usepackage[ngerman]{babel}
\usepackage{amsmath,amssymb}
\usepackage{geometry}
\geometry{margin=2.5cm}
\usepackage{tcolorbox}
\tcbuselibrary{breakable}

\newtcolorbox{merkbox}[1][]{colback=blue!5, colframe=blue!40, fonttitle=\bfseries, title={Merke}, breakable, #1}
\newtcolorbox{analogie}[1][]{colback=green!5, colframe=green!40, fonttitle=\bfseries, title={Analogie}, breakable, #1}

\title{Cheat-Sheet -- Untervektorr\"aume}
\author{}
\date{}

\begin{document}
\maketitle

%-----------------------------------------------------------------------------
\section*{1. Was ist ein Vektorraum?}

Ein \textbf{Vektorraum} ist einfach eine Menge von ``Objekten'', mit denen man zwei Dinge machen kann:
\begin{itemize}
    \item \textbf{Addieren:} Zwei Objekte zusammenz\"ahlen $\rightarrow$ Ergebnis ist wieder so ein Objekt.
    \item \textbf{Skalarmultiplikation:} Ein Objekt mit einer reellen Zahl multiplizieren $\rightarrow$ Ergebnis ist wieder so ein Objekt.
\end{itemize}

\begin{analogie}
Stell dir einen \textbf{Club} vor. Der Club hat Regeln (die Vektorraum-Axiome). Jedes Mitglied ist ein ``Vektor''. Die Regeln sagen zum Beispiel: Wenn du zwei Mitglieder ``addierst'', kommt wieder ein Mitglied raus. Wenn du ein Mitglied ``skalierst'', kommt wieder ein Mitglied raus.
\end{analogie}

\textbf{Beispiele f\"ur Vektorr\"aume:}
\begin{itemize}
    \item $\mathbb{R}^n$ -- die ``normalen'' Vektoren, z.B.\ $(1, 2, 3) \in \mathbb{R}^3$
    \item $\mathbb{R}^{n \times n}$ -- alle $n \times n$-Matrizen
    \item $\mathbb{R}[x]$ -- alle Polynome mit reellen Koeffizienten (z.B.\ $3x^2 + x - 5$)
    \item $\mathrm{FUNC}(\mathbb{R})$ -- alle Funktionen $f \colon \mathbb{R} \to \mathbb{R}$
\end{itemize}

Ein Vektorraum hat viele Regeln (Axiome): Assoziativit\"at, Kommutativit\"at, Nullelement, inverse Elemente, Distributivgesetze usw.\ -- insgesamt 8+ St\"uck. Die alle zu pr\"ufen w\"are m\"uhsam. Deshalb gibt es einen Trick...

%-----------------------------------------------------------------------------
\section*{2. Was ist ein Untervektorraum (Unterraum)?}

Ein \textbf{Untervektorraum} (kurz: Unterraum, abgek\"urzt UR) einer Menge $V$ ist eine \textbf{Teilmenge} $W \subseteq V$, die \textbf{selbst wieder ein Vektorraum} ist.

\begin{analogie}
Denk an eine \textbf{Firma} (= der gro\ss e Vektorraum $V$). Eine \textbf{Abteilung} innerhalb der Firma (= der Unterraum $W$) ist ein Teil der Firma. Aber die Abteilung muss alle Firmenregeln einhalten -- sie darf keine eigenen Regeln machen, die gegen die Firmenregeln versto\ss en.

Also: $W$ ist eine Teilmenge von $V$, und $W$ muss sich an alle Regeln von $V$ halten.
\end{analogie}

%-----------------------------------------------------------------------------
\section*{3. Der Unterraum-Test -- die 3 magischen Bedingungen}

\begin{merkbox}[title={Der Unterraum-Test (URT)}]
Sei $V$ ein Vektorraum und $W \subseteq V$. Dann ist $W$ ein Unterraum von $V$ genau dann, wenn:

\medskip
\textbf{(U1) Nullvektor:} $\vec{0} \in W$ \quad (``Der leere/neutrale Zustand muss dabei sein'')

\medskip
\textbf{(U2) Abgeschlossen unter Addition:}\\
F\"ur alle $u, v \in W$ gilt: $u + v \in W$\\
(``Wenn ich zwei Dinge aus $W$ addiere, muss das Ergebnis wieder in $W$ sein'')

\medskip
\textbf{(U3) Abgeschlossen unter Skalarmultiplikation:}\\
F\"ur alle $u \in W$ und alle $\lambda \in \mathbb{R}$ gilt: $\lambda \cdot u \in W$\\
(``Wenn ich etwas aus $W$ mit irgendeiner Zahl multipliziere, muss das Ergebnis wieder in $W$ sein'')
\end{merkbox}

\textbf{Warum reichen diese 3 Bedingungen?} Normalerweise m\"usste man ALLE 8+ Axiome eines Vektorraums pr\"ufen. Aber da $W$ eine Teilmenge von $V$ ist und $V$ schon ein Vektorraum ist, ``erbt'' $W$ automatisch die meisten Regeln von $V$:
\begin{itemize}
    \item Assoziativit\"at? Gilt schon in $V$ f\"ur alle Elemente, also auch f\"ur die in $W$.
    \item Kommutativit\"at? Gilt schon in $V$, also auch in $W$.
    \item Distributivgesetze? Genauso.
\end{itemize}

Die einzigen Dinge, die man EXTRA pr\"ufen muss, sind:
\begin{itemize}
    \item Ist das Nullelement dabei? (U1)
    \item Bleibt man ``drinnen'', wenn man addiert? (U2)
    \item Bleibt man ``drinnen'', wenn man skaliert? (U3)
\end{itemize}

\begin{analogie}
Die Firma hat schon alle Gesetze. F\"ur eine Abteilung musst du nur pr\"ufen: (1) Hat die Abteilung einen Chef (Nullelement/neutrales Element)? (2) Wenn zwei Leute in der Abteilung zusammenarbeiten, bleibt das Ergebnis in der Abteilung? (3) Wenn jemand aus der Abteilung eine Aufgabe skaliert (vergr\"o\ss ert/verkleinert), bleibt das in der Abteilung?
\end{analogie}

%-----------------------------------------------------------------------------
\section*{4. Strategie}

\begin{merkbox}[title={Strategie-\"Ubersicht}]
\textbf{Wenn du zeigen willst, dass $W$ ein Unterraum IST:}\\
$\rightarrow$ Beweise ALLE 3 Bedingungen (U1), (U2), (U3). Jede einzelne.

\medskip
\textbf{Wenn du zeigen willst, dass $W$ KEIN Unterraum ist:}\\
$\rightarrow$ Finde EIN konkretes Gegenbeispiel, das EINE der 3 Bedingungen verletzt.\\
(Tipp: Am einfachsten ist oft (U1) -- pr\"ufe, ob der Nullvektor drin ist!)
\end{merkbox}

%-----------------------------------------------------------------------------
\section*{5. Typische Fallen -- Woran erkenne ich sofort KEIN Unterraum?}

\begin{merkbox}[title={Rote Flaggen -- KEIN Unterraum}]
\textbf{Falle 1: Inhomogene Gleichungen}\\
Wenn die Bedingung so aussieht wie $a + 2c = b - 3$ (also eine Konstante $\neq 0$ auf einer Seite steht), dann ist der \textbf{Nullvektor meistens nicht drin}. Denn: Setzt man alle Variablen $= 0$, wird die Gleichung zu $0 = -3$, was falsch ist.

\medskip
Vergleiche: $a + 2c = b - 3d$ ist \textbf{homogen} (alle Terme enthalten Variablen, keine ``nackte'' Konstante). Setzt man alles $= 0$: $0 = 0$ \checkmark

\bigskip
\textbf{Falle 2: Ungleichungen ($\leq$, $<$, $\geq$, $>$)}\\
Meistens \textbf{nicht abgeschlossen unter Skalarmultiplikation}. Warum? Wenn du mit $-1$ multiplizierst, dreht sich das Vorzeichen um, und $\leq$ wird zu $\geq$!

\bigskip
\textbf{Falle 3: Produkte, Quadrate, Betr\"age}\\
Bedingungen wie $ac \leq 0$, $A^2 = 0$, $|a| \leq |c|$, $p(1) \cdot p(2) = 0$ sind meistens \textbf{nicht abgeschlossen unter Addition}. Nichtlineare Bedingungen vertragen sich schlecht mit Addition.

\bigskip
\textbf{Falle 4: $\neq$ (Ungleich)}\\
Wenn die Bedingung ein $\neq$ enth\"alt, ist der Nullvektor oft nicht drin, oder die Abgeschlossenheit scheitert.
\end{merkbox}

\begin{merkbox}[title={Gr\"une Flagge -- JA, Unterraum!}]
\textbf{Homogene lineare Bedingungen sind IMMER Untervektorr\"aume.}\\
Das hei\ss t: Wenn die Bedingung eine lineare Gleichung ist, bei der auf einer Seite $0$ steht (oder sich alles zu $0$ vereinfachen l\"asst), dann ist es ein Unterraum.

\medskip
Beispiele: $a + 2c = b - 3d$ (\"aquivalent zu $a - b + 2c + 3d = 0$), $A^\top = -A$, $p(1) = 0$, $f(-x) = -f(x)$.
\end{merkbox}

\end{document}
